\subsubsection{One-Lane Gather and Scatter Steps}

(:ref:VECTOR.instructions.VGATHER1:) and (:ref:VECTOR.instructions.VSCATTER1:) are one-lane vector memory
instructions whose repeat eligibility is limited to unconditional \texttt{REP}.
Their canonical operand order is

\begin{verbatim}
VGATHER1.T  Pn, [address], Vd
VSCATTER1.T Pn, Vs, [address]
\end{verbatim}

The read-write lane cursor \texttt{Ri} appears inside every address expression.
The vector address operand \texttt{Va} has no suffix; the instruction suffix
selects both the transferred element width and the width of \texttt{Va[Ri]}.
The encoded widths are \texttt{B/W/L/Q}.  The \texttt{H/S/D} spellings are
width-only aliases for \texttt{W/L/Q}, preserve raw element bits, require no
floating-point extension, and disassemble canonically as \texttt{W/L/Q}.

For element width $E$ bytes, let $i=unsigned(Ri)$ and $N=VLEN/E$.  The range
check precedes predicate evaluation.  If $i\mathrel{\ge}N$, the instruction
raises (:ref:VECTOR.events.EXCEPTION.VECTOR_LANE_INDEX_OUT_OF_RANGE:) and commits nothing.
Otherwise predicate bit \texttt{Pn[iE]} governs the step.  A false bit issues
no memory transaction, increments \texttt{Ri} modulo $2^{64}$, and commits.  A
true bit performs exactly one $E$-byte access through the default data segment.
On success, gather replaces only lane $i$ of \texttt{Vd}, or scatter stores
lane $i$ of \texttt{Vs}; the instruction then increments \texttt{Ri} and
commits.  A memory fault leaves \texttt{Ri}, vector state, and memory unchanged.
Both instructions preserve \texttt{FLAGS}.

The address forms, before default-data-segment translation, are:

\begin{center}
\begin{tabular}{ll}
\toprule
Form & Offset \\
\midrule
scalar stride & $Rb + unsigned(Ri)\mathbin{\cdot}signed64(Rs)$ \\
absolute 32 & $zero\_extend_{64}(Va.L[i])$ \\
absolute 64 & $Va.Q[i]$ \\
vector byte displacement & $Rb + sign\_extend_{64}(Va.E[i])$ \\
vector index & $Rb + zero\_extend_{64}(Va.E[i])\,E$ \\
indexed plus \texttt{disp8s/16s/32s} & preceding indexed offset plus the signed (:term:base.terms.payload:) \\
indexed plus \texttt{disp64} & preceding indexed offset plus the exact 64-bit (:term:base.terms.payload:) \\
\bottomrule
\end{tabular}
\end{center}

Address arithmetic is modulo $2^{64}$.  The absolute forms are offsets in the
default data segment; they do not bypass segmentation.  Absolute 32 is the
fixed \texttt{L/S} width and zero-extends \texttt{Va.L[i]}; absolute 64 is the
fixed \texttt{Q/D} width. In indexed syntax the element width $E$ directly
supplies the scale. The signed 8-, 16-, and 32-bit
displacements and exact 64-bit displacement are appended (:term:base.terms.payload|plural:).

Every base form requires \texttt{Ri != Rb}; scalar stride also requires
\texttt{Ri != Rs}, while \texttt{Rb == Rs} is permitted.  \texttt{Va == Vd}
for gather and \texttt{Va == Vs} for scatter are permitted.  The existing
vector/gather/scatter MMIO-operation prohibition applies to these one-element
transactions.
